\documentclass[10pt,a4paper]{amsart}
\usepackage[margin=19mm]{geometry}
\usepackage{amsmath,amssymb,mathtools}
\usepackage[scr=rsfs,cal=euler]{mathalfa}
\usepackage{xfrac}
\usepackage[unicode,hidelinks,bookmarks=false]{hyperref}
\newcommand{\ie}{i.e.\ }
\newcommand{\cf}{cf.\ }
\newcommand{\resp}{respectively\ }
\newcommand{\Subsection}{\subsection}
\providecommand{\nobreakdash}{\nobreak\hbox{-}}
\setlength{\emergencystretch}{2em}
\multlinegap0pt
\newtheorem{theorem}{Theorem}
\newcommand{\AR}{\operatorname{AR(1)}}
\newcommand{\Pm}{{\mathbb P}}
\newcommand{\eps}{{\varepsilon}}
\newcommand{\nat}{{\mathbb N}}
\newcommand{\rr}{{\mathbb R}}
\title{Max-semistable extremal behavior of AR(1)-processes connected with Bernoulli convolutions}
\author{Peter Kern, Alef Sterk}
\date{Source: arXiv:2608.14155v1 (CC BY 4.0)}
\begin{document}
\maketitle

\noindent\emph{Source and licence.} Max-semistable extremal behavior of AR(1)-processes connected with Bernoulli convolutions. Version arXiv:2608.14155v1 (CC BY 4.0), distributed by arXiv under the Creative Commons Attribution 4.0 International licence, http://creativecommons.org/licenses/by/4.0/, as recorded in the arXiv licence metadata for this exact version and shown on the abstract page https://arxiv.org/abs/2608.14155v1. Full licence notice: \texttt{licenses/arxiv-2608.14155-CC-BY-4.0.txt}.\par\smallskip

\noindent\textit{Editorial scope.} Counted unit: the article's theorem merge (source0.tex lines 320--325). The article's complete original proof of it is delivered with it (source0.tex lines 329--363, one \texttt{proof} environment of 35 lines, which the article places away from the statement). No mathematics is written, repaired or re-derived: the statement and the proof are the article's own lines, in the article's own notation, and no retained line is substituted or re-flowed.  Retained uncounted as the source-local context the proof needs: lines 235--239.  Nothing was omitted from any retained span, so the package carries no omission record.  No retained line contains a citation, so the wrapper carries no bibliography and no bibliography entry.  No retained line contains a figure environment, an image-inclusion command or drawing code, so no image is delivered and the package declares no asset dependency.  Editorial formatting only, in addition: an AMS \texttt{amsart} preamble that restates the article's own declarations and macros used by the retained lines, namely the environments \texttt{theorem} on separate counters, together with the standard packages those lines need.  The retained environments print in this document as Theorem 1 in order of appearance, whereas the article prints the same items with the numbering of its own full document.  The complete native source of the article as distributed in the arXiv e-print for this exact version is bundled unchanged as \texttt{sources/207638/source0.tex}.
\begin{theorem}\label{maxsdoa}
For $\beta\in(\frac12,1)$ and $n\in\nat$ let $a_{n}=\beta^{-n}$, $b_{n}=1$ and $k(n)=\lfloor q^{-n}\rfloor$, then for all $x<0$ we have as $n\to\infty$
$$\Pm\left(a_{n}(M_{k(n)}^{\circ}-b_{n})\leq x\right)\to\exp\left(-(-x)^{\log q/\log\beta}\nu_{\beta,q}^{\to}(\log(-x))\right).$$
This shows that the stationary distribution $F_{\beta,p}$ of the $\AR$ belongs to the domain of geometric partial attraction of a max-semistable distribution with parameters $\alpha=\log q/\log\beta>0$ and $c=q^{-1}>1$.
\end{theorem}

\begin{theorem}[\bf Merge theorem]\label{merge}
With the above notations we get for all $\beta\in(0,1)$ as $n\to\infty$
\begin{equation}\label{mergecomp}
\sup_{x\in\rr}\left|\Pm\left(a_{n}^{\circ}(M_{n}^{\circ}-1)\leq x\right)-H^{r_{n}}_{\beta,q}(x)\right|\to0.
\end{equation}
\end{theorem}

\begin{proof}[Proof of Theorem \ref{merge}]
It suffices to prove
\begin{equation}\label{mergesuff}
\sup_{x\leq0}\left|\Pm\left(a_{n}^{\circ}(M_{n}^{\circ}-1)\leq x\right)-H^{r_{n}}_{\beta,q}(x)\right|\to0.
\end{equation}
\underline{Step 1:} Consider a subsequence $n'=k(m_{n'})r_{n'}$ such that $r_{n'}\to r$ for some $r\in[1,c]$. Then for any sequence $(x_{n'})\subseteq(\-\infty,0]$ with $x_{n'}\to x\leq0$ we get as in the proof of Theorem \ref{maxsdoa}
\begin{align*}
\Pm\left(a_{n'}^{\circ}(M_{n'}^{\circ}-1)\leq x_{n'}\right) & =\Pm\left(a_{m_{n'}}(M_{k(m_{n'})r_{n'}}^{\circ}-1)\leq x_{n'}\right) =\Pm(Y_{0}\leq a_{m_{n'}}^{-1}x_{n'}+1)^{k(m_{n'})r_{n'}}\\
& =F_{\beta,p}(a_{m_{n'}}^{-1}x_{n'}+1)^{k(m_{n'})r_{n'}}=\left(1-F_{\beta,q}(-a_{m_{n'}}^{-1}x_{n'})\right)^{k(m_{n'})r_{n'}}\\
& =\left(1-(-a_{m_{n'}}^{-1}x_{n'})^{\log q/\log\beta}\nu_{\beta,q}(\log(-a_{m_{n'}}^{-1}x_{n'}))\right)^{k(m_{n'})r_{n'}}\\
& =\exp\left(k(m_{n'})\log\left(1-(-a_{m_{n'}}^{-1}x_{n'})^{\log q/\log\beta}\nu_{\beta,q}(\log(-a_{m_{n'}}^{-1}x_{n'}))\right)\right)^{r_{n'}}\\
& \sim\exp\left(-k(m_{n'})(-a_{m_{n'}}^{-1}x_{n'})^{\log q/\log\beta}\nu_{\beta,q}^{\to}(\log(-x_{n'}))\right)^{r_{n'}}\\
& =\exp\left(-\lfloor q^{-m_{n'}}\rfloor q^{m_{n'}}(-x_{n'})^{\log q/\log\beta}\nu_{\beta,q}^{\to}(\log(-x_{n'}))\right)^{r_{n'}}\\
& \to\exp\left(-(-x)^{\log q/\log\beta}\nu_{\beta,q}^{\to}(\log(-x))\right)^{r}=H_{\beta,q}^{r}(x),
\end{align*}
since $\nu_{\beta,q}^{\to}$ is a $\log(1/\beta)$-periodic and continuous function.\\[1ex]

\underline{Step 2:} Let $F_{n'}(x)=\Pm\left(a_{n'}^{\circ}(M_{n'}^{\circ}-1)\leq x\right)$. Since $H_{\beta,q}$ is continuous, it follows from step 1 that $F_{n'}(x)\to H_{\beta,q}^{r}(x)$ uniformly on compact subsets of $x\leq0$, whenever $(n')$ is a subsequence such that $n'=k(m_{n'})r_{n'}$ with $r_{n'}\to r\in[1,c]$. For $\eps>0$ choose $M>0$ such that for all $x<-M$ we have $F_{n'}(x)<\eps/2$ (for all $n'$ by Prohorov's theorem) and $H_{\beta,q}^{r}(x)<\eps/2$. Then we get
$$\limsup_{n'\to\infty}\sup_{x\leq0}\left|F_{n'}(x)-H_{\beta,q}^{r}(x)\right|\leq\eps+\limsup_{n'\to\infty}\sup_{-M\leq x\leq0}\left|F_{n'}(x)-H_{\beta,q}^{r}(x)\right|=\eps,$$
which shows that
\begin{equation}\label{mergesuff1}
\sup_{x\leq0}\left|F_{n'}(x)-H^{r}_{\beta,q}(x)\right|\to0
\end{equation}
as $n'=k(m_{n'})r_{n'}\to\infty$, whenever $r_{n'}\to r\in[1,c]$.\\[1ex]

\underline{Step 3:} With the same arguments as in step 2, we get
\begin{equation*}
\sup_{x\leq0}\left|H^{r_{n'}}_{\beta,q}(x)-H^{r}_{\beta,q}(x)\right|\to0,
\end{equation*}
whenever $r_{n'}\to r\in[1,c]$. Combining this with \eqref{mergesuff1} we get
\begin{equation}\label{mergesuff2}
\sup_{x\leq0}\left|F_{n'}(x)-H^{r_{n'}}_{\beta,q}(x)\right|\to0
\end{equation}
as $n'=k(m_{n'})r_{n'}\to\infty$, whenever $r_{n'}\to r\in[1,c]$. Since every subsequence of the natural numbers $n=k(m_{n})r_{n}$ contains a further subsequence $n'$ such that $r_{n'}\to r\in[1,c]$ and hence \eqref{mergesuff2} holds, we conclude that \eqref{mergesuff} is valid.
\end{proof}

\end{document}
